\chapter{Background and Related Work}\label{ch:related}
\section{Influence maximization and diffusion}
An influence-maximization objective depends on how activation spreads. The Independent Cascade model assigns transmission opportunities to newly active nodes, with edge-specific success probabilities. Expected final spread can then define a set function over seeds. The classical analysis of Kempe, Kleinberg, and Tardos establishes computational hardness and exploits monotonicity and submodularity for approximation in appropriate diffusion models \cite{kempe2003}. These results explain why the marginal benefit of a seed, rather than its degree alone, matters.

Efficient seed selection must account for overlap. Two high-degree seeds may reach nearly the same neighborhood. CELF uses lazy evaluation of marginal gains to reduce repeated work in greedy submodular selection \cite{leskovec2007}. The saved CELF result in this project is a reference denominator, not a proof that the learned model reaches a fixed fraction of a globally optimal solution. A ratio to a heuristic reference should not be confused with a fraction of all vertices.

The project deliberately uses one-step outgoing IC with transmission probability one. Its spread is the size of the union of the selected seeds and their outgoing neighbors. This is a deterministic coverage objective for a fixed graph and seed set. Multi-hop stochastic cascades raise additional issues not tested by the current evaluation. The distinction is important because a mechanism can behave differently under long-range propagation, repeated diffusion simulation, or edge probabilities below one.

\section{Learning structural influence scores}
A GNN replaces repeated combinatorial decisions with a learned map from graph structure to scores. Message passing can represent structural interactions, while a differentiable objective can encourage a useful covering set. FastCover introduces an unsupervised framework with a graph reversed attention network, GRAT, and obtains a seed set from learned node scores \cite{ni2021}. This provides architectural context for the released GRAT components inherited by the present project.

PrivIM specifically studies private GNN-based influence maximization and includes Sensitivity-Constrained Sampling and Boundary-Enhanced Sampling \cite{hou2025}. The current reproduction retains those sampling components within each stable slot. Its private mechanism, however, uses a separately documented locality argument: fixed slot membership bounds the number of affected slot contributions, each of which is clipped. The dissertation does not import a historical privacy claim merely because the same sampling function is called. The frozen accountant and adjacency definition must match the query that is actually trained.

There is also a reproduction distinction between the architecture in a mathematical appendix and a released implementation. The current runs use the legacy released GRAT layer, with biased linear maps and attention computed from pre-transformation node features. They do not use the separately available paper-appendix variant. All four final configurations retain that model identity. A different attention activation, feature construction, or parameterization would change the experiment and cannot silently be substituted when explaining results.

The thesis thus builds on existing GNN influence scoring and PrivIM components. It does not claim to introduce the first GNN influence maximizer or the first node-private influence-maximization method. Its question concerns whether a particular heterogeneous node policy, applied through a frozen training mechanism, produces a reliable utility gain against a strict uniform control.

\section{Differential privacy and private optimization}
Differential privacy bounds how much the output distribution can change between neighboring datasets. Its parameters describe a probabilistic relation between distributions rather than the amount of information an attacker happened to recover in a particular test. Postprocessing and composition provide essential design rules \cite{dwork2014}. The adjacency relation remains indispensable: a guarantee about changing one edge is different from a guarantee about adding or removing a node and all its incident edges.

DP-SGD develops private neural optimization using bounded gradient contributions, noise, and a privacy accountant \cite{abadi2016}. On ordinary records, the contribution unit is comparatively direct. A graph violates the naive analogy between one node and one independently affected gradient because a node may appear in many subgraphs and affect neighboring representations. The present mechanism therefore clips both individual position gradients and their slot sums. This is a structural sensitivity construction in addition to a numerical norm restriction.

R\'enyi differential privacy expresses privacy using divergence at orders $\alpha>1$. It supports additive composition and can be converted to $(\epsilon,\delta)$ guarantees \cite{mironov2017}. The accountant in this study uses the full 100-step Gaussian transcript with sensitivity 12. It does not assume that the 100 updates are independent model releases, nor that a nominal batch rate supplies amplification without a corresponding sampling theorem.

Privacy amplification by subsampling requires a precise sampling model and compatible adjacency. Balle, Barthe, and Gaboardi analyze amplification using couplings and divergences \cite{balle2018}. In this project, the additional personalized amplification is applied once to the complete base training mechanism, because each node's outer participation is drawn once before training. The inner RWR/SCS/BES procedure is not counted as an extra unproved amplification factor.

\section{Node-private graph neural networks}
Prior node-private GNN research demonstrates that graph dependencies need explicit treatment. Daigavane et al. analyze node-private learning with sensitivity and graph-aware amplification \cite{daigavane2021}. GAP separates encoding, private aggregation, and classification, with edge- and node-level protection analyzed for its defined pipeline \cite{sajadmanesh2023}. DPAR decouples feature projection and message passing using privately computed approximate Personalized PageRank neighborhoods \cite{zhang2024}.

These examples protect different computations from the current influence-ranking task. Their guarantees cannot be transferred by relabeling a loss or replacing a classifier with Top-50 ranking. Each method controls how protected graph information enters its queries and subsequent inference. The present stable-slot mechanism instead bounds arbitrary within-slot changes after node removal, while retaining the original subgraph-based objective inside each slot.

This comparison also distinguishes two meanings of ``personalized.'' Personalized PageRank concerns a graph-walk starting distribution; it does not inherently assign different DP budgets to people. Personalized prediction or federated model adaptation likewise need not imply personalized privacy requirements. The relevant policy in this dissertation is an explicit per-identity upper-bound vector, held fixed across neighboring graph inputs.

\section{Personalized privacy in graph learning}
There is already published work explicitly combining personalized privacy and GNNs. Yuan et al.'s PDPGNN assigns budgets using node degree and employs weighted aggregation \cite{yuan2026}. It addresses degree-dependent utility degradation. The present project instead fixes an external public Registry with target classes 1, 2, and 4 and uses one-time participation to meet those requirements. Its policy is not recomputed from the changing private degree sequence.

He et al. propose PPGNN for decentralized local graph data, combining personalized local perturbation with a weighted calibration strategy \cite{he2026}. That trust model differs from central private training over a fixed prepared graph. The paper is cited here in its verifiable arXiv form; its repository note about IEEE TKDE is not treated as independently verified final publication metadata. Its existence reinforces that the broad combination of personalization and graph privacy is not itself a first-of-its-kind claim.

The distinguishing elements of the present study are the IM task, the fixed target policy, the stable overlapping-slot construction, the full-coordinate training rule, and the complete statistical evaluation. They define a concrete research contribution only in combination with a correct scope of analysis. No theoretical or experimental comparison against PDPGNN or PPGNN is reported, because those implementations were not run under this frozen protocol. A literature comparison is not an empirical baseline result.

\section{High-dimensional optimization and stability}
A Gaussian mechanism can have a sensitivity bound independent of parameter count while injecting noise in every coordinate. If $Z\sim\mathcal{N}(0,\sigma^2 I_d)$, then $\E\|Z\|_2^2=d\sigma^2$. This elementary identity helps explain why the geometry of private optimization deserves attention. It does not imply that the same scaling holds for errors in logits or Coverage, which are nonlinear functions of the learned parameters and graph.

The current mechanism trains 2340 coordinates; a historical frozen-backbone comparator trains only 33. Their common per-coordinate scale can coexist with very different total parameter-space noise energy and clipping behavior. Historical results make that distinction worth investigating but do not isolate it causally. Different trainable subspaces, signal directions, and independent secret trajectories change together.

Stability must also be defined operationally. A small standard deviation can coexist with consistently poor utility, while a high mean can conceal severe losses in some runs. This dissertation therefore reports mean, sample SD, minima, maxima, and the frequency of Coverage below 50\%. Its confirmatory question concerns an average effect; it does not automatically test a simultaneous claim of higher mean and lower risk.

\section{Position of the present research}
The reviewed literature supplies the IM formulation, GRAT architecture, PrivIM sampling, differential privacy principles, RDP, and amplification tools. The project combines these within a fixed-domain personalized node-private training design and tests a specific utility hypothesis. Its conditional theorem concerns model training. Its empirical study concerns saved scores under a frozen evaluation protocol. Neither establishes private publication of all research artifacts.

This positioning leaves room for a substantive dissertation even when the primary utility comparison is inconclusive. The question is answered with an estimate and its uncertainty; the mechanism's limits are documented rather than hidden. Stronger novelty, universal superiority, or end-to-end privacy claims would require evidence beyond the present literature review, proof boundary, and five fixed datasets.
